\chapter{Conclusións}
\label{chap:conclusions}

\lettrine{N}{este} derradeiro capítulo recóllese a validación realizada cos expertos a partir
dunha reunión e dunha proba práctica do sistema. Ademais, resúmense as principais conclusións
do traballo, analízanse as limitacións detectadas e propóñense posibles liñas futuras de mellora.

\section{Validación con expertos}
Nos últimos días de desenvolvemento do sistema, cando xa se dispoñía da versión final da CNN,
organizouse unha reunión con expertos da Facultade de Ciencias. Máis concretamente, participaron
membros do grupo BIOCOST (Bioloxía Costeira), que cederan o material empregado durante este
traballo. A reunión contou coa participación de Sara Barrientos de la Llana, investigadora
postdoutoral, e Nerea Alvite Bergara, investigadora predoutoral, ambas pertencentes ao grupo
mencionado.

\subsection{Resumo da reunión}
A reunión comezou cunha presentación do sistema desenvolvido, na que se explicaron as métricas
empregadas na avaliación da CNN e se mostraron os resultados obtidos coa versión final. A recepción
foi positiva, unha vez explicado o efecto de arrastre que pode producir un erro en \texttt{HPI} sobre
a aplicabilidade de \texttt{IVR}, así como o significado das novas métricas condicionais, o equipo
experto indicou que os resultados eran coherentes co comportamento esperado.

Aínda así, co obxectivo de realizar unha comprobación máis obxectiva, propúxose unha proba práctica:
as dúas investigadoras avaliarían manualmente un subconxunto de 200 imaxes extraídas do
\textit{dataset}. Este subconxunto non se xerou mediante unha mostraxe completamente aleatoria, senón
mediante unha selección estratificada para que estivese representada toda a escala de \texttt{HPI} e,
dentro dos casos nos que \texttt{IVR} aplica, tamén a diversidade das notas de \texttt{IVR}.

Para iso, fixouse un mínimo aproximado do 10\% das imaxes para cada valor de \texttt{HPI} e, nos
valores \texttt{1..4}, procurouse repartir tamén as mostras entre as clases positivas de
\texttt{IVR}. Nos casos \texttt{HPI=0}, \texttt{HPI=5} e \texttt{HPI=6}, a etiqueta \texttt{IVR}
mantívose en \texttt{0}, seguindo a lóxica experta descrita anteriormente.

{\rowcolors{2}{white}{udcgray!25}
\begin{table}[htbp]
    \centering
    \footnotesize
    \setlength{\tabcolsep}{4pt}
    \begin{tabular}{p{0.12\textwidth}|p{0.14\textwidth}|p{0.62\textwidth}}
        \rowcolor{udcpink!25}
        \textbf{\texttt{HPI}} & \textbf{Imaxes} & \textbf{Distribución de \texttt{IVR} dentro do grupo} \\\hline
        0 & 34 & \texttt{IVR=0}: 34 \\
        1 & 47 & \texttt{1}: 7, \texttt{2}: 6, \texttt{3}: 5, \texttt{4}: 5, \texttt{5}: 5, \texttt{6}: 7, \texttt{7}: 12 \\
        2 & 27 & \texttt{1}: 5, \texttt{2}: 4, \texttt{3}: 3, \texttt{4}: 3, \texttt{5}: 3, \texttt{6}: 4, \texttt{7}: 5 \\
        3 & 23 & \texttt{1}: 4, \texttt{2}: 3, \texttt{3}: 3, \texttt{4}: 3, \texttt{5}: 3, \texttt{6}: 3, \texttt{7}: 4 \\
        4 & 22 & \texttt{1}: 5, \texttt{2}: 3, \texttt{3}: 3, \texttt{4}: 3, \texttt{5}: 3, \texttt{6}: 2, \texttt{7}: 3 \\
        5 & 26 & \texttt{IVR=0}: 26 \\
        6 & 21 & \texttt{IVR=0}: 21 \\
    \end{tabular}
    \caption{Distribución do subconxunto de 200 imaxes empregado para a avaliación manual por parte dos expertos.}
    \label{tab:dataset-validacion-expertos}
\end{table}
}

Durante a xeración deste subconxunto só se detectou unha limitación: para a combinación
\texttt{HPI=4, IVR=6} había dúas imaxes dispoñibles, aínda que o obxectivo mínimo eran tres. Neste
caso empregouse a mellor aproximación posible, sen alterar o tamaño total do subconxunto. Ademais,
exportáronse dous ficheiros separados: un modelo de revisión para os expertos, sen información das
etiquetas orixinais, e unha chave privada coas respostas para poder comparar posteriormente as
avaliacións.

Con esta avaliación manual, buscábase comprobar ata que punto o sistema desenvolvido se podía
comparar coa avaliación experta. Ademais, a proba permitía analizar se, incluso entre expertos, a
variabilidade da nota podía ser elevada en certos casos pola subxectividade da detección visual do
dano ou da atribución da depredación.

\subsection{Resultados da proba}

\section{Conclusións do traballo}
\subsection{Cumprimento dos obxectivos}
\subsection{Aportacións principais}

\section{Limitacións do sistema}

\section{Liñas futuras}
